\emph{Turing机}是一种理想化的计算机制。它由一条单向无限的\emph{磁带}组成，磁带被划分为方格，每个方格可以包含预先确定的字母表中的某个\emph{符号}。机器通过沿磁带移动一个\emph{读写头}来运行。它还可以处于预先确定的若干\emph{状态}之一。读写头的动作由一组指令决定；每条指令都以机器处于某个特定状态并读取某个特定符号为条件，并指定机器将在当前方格写入哪个符号、将读写头向左或向右移动一格还是保持不动，以及它将切换到哪个状态。如果磁带上包含某个特定的\emph{输入}（表示为磁带上的符号序列），并且机器被置于指定的起始状态，同时读写头读取输入最左侧的方格，那么指令集将逐步决定机器的一系列\emph{格局}：磁带内容、读写头位置以及机器状态。一旦机器遇到某个格局，其中指令集不包含针对当前读取的符号与所处状态组合的指令，机器就\emph{停机}，此时磁带的内容即为\emph{输出}。

数字可以非常容易地表示为Turing机磁带上的一列竖杠。如果存在一台Turing机，使得每当它以~$n$ 的一进制表示作为输入启动时，最终都会停机，并且其磁带上包含作为输出的~$f(n)$ 的一进制表示，我们就称函数 $\Nat \to \Nat$ 是\emph{Turing可计算的}。许多常见的算术函数都能被证明是Turing可计算的（有些容易，有些稍费周折）。人们还提出了许多Turing机以外的其他计算模型；而结果总是表明，这些模型上可计算的算术函数也是Turing可计算的。这被视为对\emph{Church-Turing论题}的支持，该论题断言：每个能有效计算的算术函数都是Turing可计算的。